\section{Sharp lattice conditioning and what it controls}\label{sec:lattice}
We record the quantitative lattice estimate needed to distinguish coefficient growth from traversal. The orthant estimate and its sharp shear construction were established in the preceding Manuscript II checkpoints~\cite{M2prior}; the qualitative Interior Ryshkov property and the underlying Diophantine method are prior literature~\cite{DSV21,OS26}.

\begin{lemma}[Orthant norm-to-margin bound]\label{lem:orthantmargin}
If $P[v]\ge1$ for every $0\ne v\in\Z_{\ge0}^s$ and $M\ge\max(1,\norm P_F)$, then
\[
 \min_{x\ge0,\ \norm x_2=1}P[x]\ge\frac1{8s(16sM)^{s-1}}.
\]
\end{lemma}
\begin{proof}
A negative value at a real nonnegative vector would persist at a rational vector and contradict the integer inequalities after clearing denominators. Thus $P$ is copositive. Choose a unit minimizer $x$ with value $\mu$, and let its positive support have size $a$. Stationarity on the support gives $P_{SS}x_S=\mu x_S$. If $a=1$, the claim follows from $P_{ii}\ge1$. Otherwise normalize a largest coordinate to one, writing $y=x/x_j$ with $\norm y_2\le\sqrt s$.

Let $Q=\lceil\sqrt{8sM}\rceil$. The pigeonhole proof of simultaneous Dirichlet approximation gives $1\le m\le Q^{a-1}$ and an integer nonnegative vector $v$ with $v_j=m$ and $|v_i-my_i|\le1/Q$ on the other support coordinates. Nonnegativity follows since the approximated numbers are nonnegative and the error is less than one. Put $e=v-my$ and use support stationarity to get
\[
 1\le P[v]\le\mu(m\norm y+\norm e)^2+M\norm e^2
 \le4s\mu Q^{2(s-1)}+1/8.
\]
Since $Q\le4\sqrt{sM}$, the stated weaker constant follows. The argument proves strict positivity rather than presupposing it.
\end{proof}

\begin{theorem}[Cone-relative conditioning, with sharp exponent]\label{thm:lattice}
Let a rational polyhedral cone $K$ of span dimension $r$ be generated by integer vectors of norm at most $R\ge1$. If $P[v]\ge1$ for every nonzero integer $v\in K$ and $M\ge\max(1,\norm P_F)$, then
\begin{equation}\label{eq:Kmargin}
 \min_{x\in K,\ \norm x_2=1}P[x]\ge
 \frac1{8r^2R^2(16r^2R^2M)^{r-1}}.
\end{equation}
For every fixed rational polyhedral cone with relative dimension $r\ge2$, the exponent $r-1$ cannot be reduced, even on relatively perfect positive-definite forms.
\end{theorem}
\begin{proof}
Every $x\in K$ is a nonnegative combination $Ba$ of $s\le r$ independent integer generators. The pullback $B\trans PB$ has nonnegative integer minimum at least one, norm at most $sR^2M$, and $\norm{Ba}^2\le sR^2\norm a^2$. Apply Lemma~\ref{lem:orthantmargin} and enlarge constants from $s$ to $r$ to get \eqref{eq:Kmargin}.

For sharpness, use the relative lattice $\lin K\cap\Z^n$. It has a basis $B$ lying in $K$: choose a primitive relative-interior lattice vector, extend it to a lattice basis, and add sufficiently large integer multiples of it to the other basis vectors. The relative-interior vector can be made primitive by dividing its lattice gcd. Let $T_r$ have diagonal one and adjacent off-diagonal $-1/2$, and let $V_q=I-qS$ for the upper shift $S$ and integer $q\ge2$. Define the form by
\[
 P_q[Bz]=z\trans V_q\trans T_rV_qz.
\]
For every nonzero integer $z$,
\[
 T_r[z]=\tfrac12\left(z_1^2+z_r^2+\sum_{i=1}^{r-1}(z_i-z_{i+1})^2\right)\ge1.
\]
The positive interval indicators attain one and their squares span $\Ssym^r$. Since $V_q^{-1}=I+qS+\cdots+q^{r-1}S^{r-1}$ is nonnegative integral, their inverse images under $V_q$, then under $B$, lie in $K$ and prove relative perfection. The norm of $P_q$ is $O_K(q^2)$. The minimal vector $BV_q^{-1}e_r$ has norm $\Omega_K(q^{r-1})$ and value one, so the margin is $O_K(q^{-2(r-1)})$. No smaller norm exponent is possible.
\end{proof}

The bound limits minimal vector lengths. In full dimension, writing its reciprocal as $U$, every active integer vector has squared norm at most $U$. In upper-triangular coordinates its integer constraint normal has entries at most $2U$. Cramer's rule bounds vertex entries and primitive tangent-ray entries by $d![2\lceil U\rceil]^d$. Thus coefficient heights grow only logarithmically with $M$ for fixed dimension. These bounds do not count how many normal-fan cells are visited. The next theorem proves that this missing resource is genuinely independent of the arithmetic carrier.

\section{An exponential local walk inside one arithmetic face}\label{sec:paths}
Fix $r\ge2$, put $n=r+2$, and consider
\begin{equation}\label{eq:pathcone}
 K=\{(t,z_1,\ldots,z_r,y):t,y\ge0,\ |z_i|\le t\}.
\end{equation}
This is the product of a homogenized cube and a ray. It has $2^r+1>n$ extreme rays and is not simplicial. Its generalized Ryshkov set consists of $P$ with $P[v]\ge1$ for all nonzero $v\in K\cap\Z^n$. At a perfect vertex let the shadow rule use the fixed anchor $C_0$ and minimize $\ip{C_0}{R}/(-\ip A R)$ among objective-descending extreme tangent rays. We count exact adjacent moves, not arbitrary preprocessing or nonlocal factor recognition.

\begin{theorem}[Depth-zero local/global separation]\label{thm:path}
For the fixed cone \eqref{eq:pathcone} and every integer $N\ge2$, there is a rational interior target $A_N$ with the following properties.
\begin{enumerate}[label=(\roman*)]
\item Its input length is $O_n(\log N)$. It has at most $6\cdot3^r$ explicit rational Gram columns, each of height $O_n(\log N)$.
\item Its arithmetic carrier is the whole $K$, and its arithmetic reduction depth is zero. For full observation, $\beta(K,\mathrm{id})=1/\sqrt n$, independent of $N$.
\item From a fixed perfect vertex and its full-minimum anchor, the specified shadow rule has a unique choice at every crossing and makes exactly $N-1$ adjacent pivots.
\item After rational-square normalization, $\norm{A_N}_F\le1$ and its interior distance is at least
\begin{equation}\label{eq:pathradius}
 \frac{2^{-(2r+4)}}{n^2d!(N+1)^{2d}}.
\end{equation}
Every visited vertex, ray and compressed certificate has height $O_n(\log N)$.
\end{enumerate}
A fixed integral unimodular congruence makes all target entries strictly positive, preserving the local path and the conclusions up to fixed-dimensional constants.
\end{theorem}

\begin{proof}
For $k\ge1$ define a quadratic form, independent of $z$, by
\begin{equation}\label{eq:Pk}
 P_k=t^2-(2k+1)ty+(k^2+k+1)y^2.
\end{equation}
With $a=t-ky$ this is $a^2-ay+y^2$. Its allowed nonzero integer minima are exactly the vectors with
\begin{equation}\label{eq:minima}
 (t,y)=(1,0),(k,1),(k+1,1),\qquad z\in\Z^r,\quad |z_i|\le t.
\end{equation}
Indeed the binary integer form has minimum one; if $y\ge2$, completing the square gives value at least $3y^2/4\ge3$. The cases $y=0,1$ give precisely the stated possibilities. If $t=y=0$, cone membership forces $z=0$.

These minima prove perfection. A quadratic form vanishing on $(1,z,0)$ for all $z\in\{-1,0,1\}^r$ has zero $tt,tz,zz$ coefficients by interpolation on three nodes in each coordinate. Vanishing on $(k,z,1)$ next makes every $zy$ coefficient zero and imposes $2kQ_{ty}+Q_{yy}=0$. The family with $k+1$ gives $Q_{ty}=Q_{yy}=0$. Thus the active normals span $\Ssym^n$.

For $s\ge2$ the common active families at $P_{s-1},P_s$ have base coordinates $(1,0),(s,1)$. The same calculation leaves the one-dimensional annihilator
\[
 R_{s-1}=P_s-P_{s-1},\qquad (R_{s-1})_{ty}=-1,\quad(R_{s-1})_{yy}=2s.
\]
The segment between the two feasible vertices is an edge. Its length in this normalization is exactly one: a vector with base $(s+1,1)$ becomes infeasible immediately beyond $P_s$. Relative strict copositivity of $P_k$ on $K$ follows from the positive-definite binary form and $\norm z^2\le r t^2$, so the local normal cones are generated by the finitely many active vectors.

At $P_1$ the anchor obtained by summing \emph{all} minima has binary entries
\begin{equation}\label{eq:anchor}
 a=2\cdot3^r+4\cdot5^r,\quad b=3^r+2\cdot5^r,\quad c=3^r+5^r
\end{equation}
and each $z$ diagonal equal to $\zeta=4\cdot3^{r-1}+10\cdot5^{r-1}$. All other $z$ entries vanish. Note $a=2b$, $2c-b=3^r$, and $0<\zeta<a$.

Write $T=N+1/2$ and $U=N^2+N+1$. Define the unnormalized target by its binary block
\begin{equation}\label{eq:target}
 A_N|_{t,y}=\begin{pmatrix}U&T\\T&1\end{pmatrix},\qquad
 A_N|_z=(U/2)I_r,
\end{equation}
with all remaining entries zero. Its binary determinant is $3/4$, hence it is positive definite. An explicit permitted factorization starts with base vectors $(1,0),(N,1),(N+1,1)$, each of weight $1/2$, and independently assigns $z_i=-t,0,t$ with probabilities $1/4,1/2,1/4$. The $3\cdot3^r$ resulting weighted atoms have Gram sum \eqref{eq:target}. All weights are positive dyadic numbers and their normals span $\Ssym^n$ by the preceding interpolation argument. Thus $A_N\in\operatorname{int}C(K)$. Each dyadic weight is one rational square or twice one rational square, giving at most $6\cdot3^r$ Gram columns.

For $2\le s\le N$ set
\begin{equation}\label{eq:lambda}
 \lambda_s=\frac{sc-b}{T-s},\qquad C_s=C_0+\lambda_s A_N.
\end{equation}
The parameters are positive and strictly increasing, because $2c-b>0$ and $cT-b>0$. The binary part of $C_s$ has the exact decomposition
\begin{equation}\label{eq:binarynormal}
 \chi_s(1,0)(1,0)\trans+(c+\lambda_s)(s,1)(s,1)\trans,
\end{equation}
where
\begin{align*}
 \chi_s&=a-sb+\lambda_s(U-sT)\\
 &=a-bT+s(cT-b)+\tfrac34\lambda_s>0.
\end{align*}
For positivity, $s\ge2$ and $cT-b>0$ make the last expression at least $(2c-b)T+\tfrac34\lambda_s$.

Let
\[
 v_s=\frac{\zeta+\lambda_s U/2}{a+\lambda_s U}\in(0,1).
\]
Lift each base vector in \eqref{eq:binarynormal} independently by $z_i=-t,0,t$ with probabilities $v_s/2,1-v_s,v_s/2$. This gives exactly $C_s$, with strictly positive weights on normals spanning the common active hyperplane. Hence $C_s$ lies in the relative interior of the shared normal-cone facet. Its sole zero tangent ray is $R_{s-1}$, and
\[
 \ip{A_N}{R_{s-1}}=-2(T-s)<0.
\]
Between two successive crossings the positive combination of the two facet certificates uses the three base families $(1,0),(s,1),(s+1,1)$, whose normals span the full matrix space. Therefore the intervening objective is in the interior normal cone at $P_s$. Initially $C_0$ is interior at $P_1$. Finally $A_N$ itself is interior in the normal cone at $P_N$, so every objective after the last crossing remains there. This proves the exact unique shadow trajectory $P_1,P_2,\ldots,P_N$ and its $N-1$ pivots. No chosen ray can be a copositive separator: each has a negative permitted value, and the target is completely positive. The statement does not prohibit a different algorithm from recognizing the displayed factorization immediately.

For the interior radius, select $d$ independent integer atoms from the displayed target representation. Their upper-triangle coordinate matrix has integer determinant of absolute value at least one and entries bounded by $(N+1)^2$. Its inverse entries have magnitude at most $(d-1)!(N+1)^{2(d-1)}$. Every target weight is at least $w_0=2^{-(2r+1)}$. A Frobenius perturbation of norm at most
\[
 \eta=\frac{w_0}{2d!(N+1)^{2(d-1)}}
\]
can consequently be absorbed into those $d$ coefficients while keeping them positive. Normalize by the rational square $[2n(N+1)]^2$. This bounds the Frobenius norm by one and gives \eqref{eq:pathradius}. All expressions involve rational numbers with $O_n(\log N)$ bits, including $s\le N$ in the compressed normal certificates.

The carrier is $K$ because the target is an interior rational point; Theorem~\ref{thm:main} or the reserve argument applies with no reduction. For full observation the trace-one PSD norm is at least $1/\sqrt n$. Moreover $I\in C(K)$: independent signs in $(1,z,0)$ generate $\diag(1,I_r,0)$, and the $y$ axis supplies the last diagonal. Hence $\beta(K,\mathrm{id})=1/\sqrt n$.

Finally apply the fixed integral map
\[
 (t,z,y)\longmapsto(t,z+2t\mathbf1,y+t).
\]
Its determinant is one. It sends $K$ to the nonsimplicial cone $t\ge0$, $t\le z_i'\le3t$, $y'\ge t$. The congruenced target is positive definite and every entry is strictly positive. Integral minima, full-minimum anchors, trace pairings, extreme rays and the minimum-ratio rule are carried bijectively by the inverse-transpose action on forms. Norm and interior-radius estimates change only by constants depending on $n$. This finishes the proof.
\end{proof}

\begin{corollary}[No polynomial walk bound from genuine arithmetic faces]\label{cor:barrier}
For the specified fixed-start rule, no polynomial in the binary input length, rational factor height and width, binary interior condition, and arithmetic reduction depth bounds the number of adjacent pivots, even after fixing the vector cone, its arithmetic carrier, and the observation coercivity. The entire shadow objective line can remain in the same genuine interior matrix-cone face while crossing exponentially many perfect-form cells.
\end{corollary}
\begin{proof}
Take $N=2^B$ in Theorem~\ref{thm:path}. Every listed binary resource is $O_n(B)$ or fixed, the depth is zero, and the pivot count is $2^B-1$. The anchor and target are interior, hence so is their positive objective combination throughout the path.
\end{proof}

This is an obstruction to a particular proposed complexity inference, not a claim that no imaginable invariant can encode path length, and not a lower bound for every descending local rule. The earlier Manuscript II results provide complementary nonmember instances~\cite{M2prior}; they are not needed for this interior argument. The shear/perfectness architecture has antecedents in Dannenberg--Sch\"urmann~\cite{DS23}. What is isolated here is the fixed nonsimplicial, depth-zero, uniquely forced path with an explicit short rational certificate and inverse-polynomial interior radius.

